Con \(r=0.6\), el primer paso es \((-2,6)\,\mathrm K\). El mayor paso que conserva no negatividad para todos los datos no negativos es \(\Delta t=0.5\,\mathrm s\). En la malla fija, los dos multiplicadores con el paso de \(0.6\,\mathrm s\) son \(0.4\) y \(-0.8\), por lo que toda perturbación inicial decae aunque pueda aparecer una temperatura negativa. Ninguna de esas garantías cuantifica por sí sola el error respecto de la solución continua.
